\documentclass[12pt,a4paper]{article}
\usepackage[utf8]{inputenc}
\usepackage[T1]{fontenc}
\usepackage[brazil]{babel}
\usepackage[margin=2.5cm]{geometry}
\usepackage{booktabs}
\usepackage{amsmath,amssymb}
\usepackage{caption}
\usepackage{graphicx}
\usepackage{setspace}
\usepackage{indentfirst}
\usepackage{fancyhdr}
\usepackage[hidelinks]{hyperref}
\usepackage{enumitem}
\graphicspath{{docs/}}

\captionsetup{font=small,labelsep=colon}
\onehalfspacing
\emergencystretch=2em

\pagestyle{fancy}
\fancyhf{}
\fancyhead[L]{\footnotesize Determinantes e previsibilidade dos pesos de ações em fundos}
\fancyhead[R]{\footnotesize J. P. Zangrandi}
\fancyfoot[C]{\thepage}
\renewcommand{\headrulewidth}{0.4pt}

\begin{document}

%==================== CAPA ====================
\thispagestyle{empty}
\begin{center}
{\large FUNDAÇÃO GETÚLIO VARGAS}\\[2pt]
{\large ESCOLA DE ECONOMIA DE SÃO PAULO}\\[16pt]
{\normalsize Graduação em Economia}

\vspace{3.0cm}

{\large João Paulo Zangrandi}

\vspace{2.5cm}

{\Large Estudo sobre a Demanda Institucional por Ativos Financeiros}\\[6pt]
{\normalsize Versão I --- correções de inferência e extensão multiativo}

\vspace{2.5cm}

\begin{flushright}
\begin{minipage}{0.62\textwidth}
\small\setstretch{1.15}
Monografia I apresentada à Escola de Economia de São Paulo da Fundação Getúlio
Vargas como requisito da disciplina Monografia I.
\\[10pt]
Orientador: Prof.\ Maurício Ferraresi Junior.
\end{minipage}
\end{flushright}

\vfill

{\normalsize São Paulo}\\
{\normalsize 2026}
\end{center}

\newpage

%==================== RESUMO / ABSTRACT ====================
\thispagestyle{fancy}
\begin{center}{\large Resumo}\end{center}
\noindent
Este trabalho estuda os determinantes da demanda dos fundos do Itaú entre 2016 e 2021 e a
capacidade desses determinantes de prever as escolhas futuras desses fundos. O foco inicial
foi a demanda pela ação da Vale no mercado brasileiro. A abordagem segue a lógica de um
sistema de demanda institucional: a cada mês, o peso é regredido por mínimos quadrados sobre
características de cada fundo, e os coeficientes são lidos como um fator latente que varia no
tempo. O peso é explicado pelo tamanho do fundo, pelo número de cotistas e pela classificação
como fundo de cotas, com sinais estáveis nos 72 meses e que generalizam, na mesma direção,
para 92 a 99\% de um universo ampliado de quase 900 ações mantidas por fundos do Itaú ---
embora o efeito de tamanho da própria VALE3 esteja no percentil 1 dessa distribuição, um caso
extremo, não típico. O trabalho estima a velocidade de ajuste da carteira ao peso desejado por seis especificações de
mínimos quadrados, variável instrumental e efeito-fixo de fundo, com sinal positivo e
significância robustos em todas elas, ainda que a magnitude varie por um fator de dez conforme o
estimador; uma sétima família, o painel dinâmico de Arellano--Bond, não confirma essa robustez
--- o coeficiente troca de sinal conforme a forma de construir os instrumentos, um sintoma de que
o estimador é mal identificado neste desenho de $N$ pequeno e $T$ grande, e por isso não é lido
como evidência adicional. Um teste formal de Diebold--Mariano mostra que o ganho de previsão fora
da amostra do modelo de ajuste parcial sobre a previsão ingênua só é estatisticamente detectável
no horizonte de um mês, não no horizonte de três meses relevante sob a opacidade de divulgação
das carteiras. Esta versão corrige inconsistências de escala e de inferência identificadas em
revisão crítica independente, adota a captação winsorizada como especificação principal, e
acrescenta como controles o beta do próprio fundo e os fatores de risco do NEFIN-USP.

\smallskip
\noindent Palavras-chave: sistema de demanda; fundos de investimento; pesos de carteira;
fator latente dinâmico; previsão; painel dinâmico; mercado acionário brasileiro.

\bigskip
\begin{center}{\large Abstract}\end{center}
\noindent
This paper studies the determinants of demand of Itaú funds between 2016 and 2021 and the
capacity of these determinants to predict the future choices of these funds, with an initial
focus on the demand for Vale's stock. The weight is explained by fund size, number of
shareholders, and fund-of-funds classification, with signs that are stable across 72 months and
that generalize, in the same direction, to 92--99\% of an expanded universe of almost 900 stocks
held by Itaú funds --- although Vale's own size effect sits at the 1st percentile of that
cross-asset distribution, an extreme rather than typical case. The paper estimates the
portfolio's adjustment speed toward the desired weight with six OLS, instrumental-variable and
fund-fixed-effects specifications, all agreeing on a positive and significant reversion despite a
tenfold range in magnitude; a seventh family, Arellano--Bond dynamic panel GMM, does not confirm
that robustness --- the coefficient flips sign depending on how instruments are built, consistent
with weak identification in this small-$N$, large-$T$ design, and is therefore not read as
additional evidence. A formal Diebold--Mariano test shows that the out-of-sample forecast gain of
the partial-adjustment model over the naive prediction is only statistically detectable at the
one-month horizon, not at the three-month horizon relevant under portfolio-disclosure opacity.
This version corrects scale and inference issues found in an independent critical review, adopts
winsorized net flow as the main specification, and adds the fund's own beta and NEFIN-USP risk
factors as controls.

\smallskip
\noindent Keywords: demand system; mutual funds; portfolio weights; time-varying latent factor;
forecasting; dynamic panel; Brazilian equity market.

\newpage
\thispagestyle{fancy}
\tableofcontents
\newpage

%==================== NOTA DE VERSÃO ====================
\section*{Nota sobre esta versão}
\addcontentsline{toc}{section}{Nota sobre esta versão}
Este documento (``tcc I'') substitui as versões anteriores (\texttt{tcc}, \texttt{tcc\_new},
\texttt{tcc\_novo}, \texttt{novo}, \texttt{resultados\_2}) como a versão corrente do trabalho.
Ele incorpora: (i) todas as correções apontadas em uma revisão crítica independente
(\texttt{correções\_fable.pdf}, 06/07/2026) --- teste de Diebold--Mariano para a comparação de
previsão, correção do erro-padrão da velocidade de ajuste por efeito de tempo, especificação
com efeito-fixo de fundo reapresentada como co-principal, captação winsorizada como
especificação principal, efeito parcial médio corrigido para a variável binária, e sensibilidade
do erro-padrão de Newey--West à defasagem; e (ii) uma nova rodada de anotações do orientador
(09/07/2026): estimação da velocidade de ajuste por Arellano--Bond, beta do próprio fundo como
controle, fatores de risco do NEFIN-USP, e a generalização da estimação de demanda para o
universo amplo de ações mantidas pelos fundos do Itaú. A elasticidade-preço no estilo
Koijen--Yogo (o canal preço-demanda ``de volta'') fica registrada como agenda futura, não
estimada aqui --- ver Seção~\ref{sec:agenda}.

%==================== 1. INTRODUÇÃO ====================
\section{Introdução}

A composição das carteiras dos fundos de investimento é uma das informações mais ricas e ao
mesmo tempo menos exploradas do mercado de capitais brasileiro. Quando um fundo decide quanto
de cada ação carregar, ele revela, ainda que parcialmente, a sua demanda por aquele papel.
Olhar para todos os fundos de uma mesma gestora ao mesmo tempo equivale, então, a observar um
grande conjunto de decisões de demanda tomadas sob informações e incentivos parecidos --- uma
hipótese de trabalho que vale explicitar: as três entidades do grupo (Itaú Asset Management,
Itaú DTVM e Itaú Unibanco) operam sob a mesma marca e, presumivelmente, compartilham parte da
pesquisa e da governança de risco, mas não são a mesma mesa de gestão, e o grau exato de
similaridade de incentivos entre elas não é testado diretamente neste trabalho. Isso torna
possível investigar de forma organizada duas perguntas. A primeira é quais características de
um fundo explicam o peso que ele atribui a uma ação. A segunda é se esse peso pode ser previsto
de um mês para o outro.

Para tornar o problema tratável, a maior parte deste trabalho concentra-se em um caso concreto:
o peso da ação VALE3 nas carteiras dos fundos das três gestoras do grupo Itaú no período de 2016
a 2021. A escolha do peso como variável de interesse, em vez do valor financeiro ou do número de
ações, é deliberada. O peso é limitado entre zero e um, o que o torna mais simples de comparar;
é comparável entre fundos de tamanhos muito distintos; e corresponde diretamente à noção de
demanda relativa que organiza os modelos de demanda por ativos. Esta versão do trabalho amplia
adicionalmente a estimação da parte transversal do modelo (Seção~\ref{sec:multi}) para o
universo inteiro de ações que os fundos do Itaú carregam --- quase 900 papéis com amostra
suficiente ---, o que permite perguntar se o padrão encontrado para VALE3 é representativo do
que acontece com uma ação qualquer, ou se é uma particularidade daquele papel.

O arcabouço teórico usado para este trabalho é o sistema de demanda de Koijen e Yogo (2019),
que propõem tratar os preços dos ativos como resultado das demandas de investidores que
escolhem suas carteiras em função das características dos ativos. A principal vantagem dessa
abordagem é reduzir um problema de dimensão muito alta, com milhares de ativos e investidores,
a um número pequeno de parâmetros associados às características.

É necessário, porém, ser preciso quanto ao ponto em que este trabalho se afasta daquele
modelo. Em Koijen e Yogo, a demanda do investidor $i$ pelo ativo $n$ é função das
características do ativo, e o vetor de sensibilidades varia por investidor; a identificação
enfrenta a endogeneidade entre preço e demanda por meio de um instrumento construído a partir
das características dos demais investidores. Aqui a lógica é, em sentido estrito, invertida em
três aspectos.

Primeiro, na análise principal o ativo é fixo, VALE3, e o que varia entre observações são as
características do fundo (a Seção~\ref{sec:multi} relaxa esse ponto, estimando o mesmo tipo de
regressão para muitos ativos simultaneamente). Segundo, a sensibilidade das escolhas em relação
às características de cada fundo é comum a todos os fundos em cada mês, em vez de específica ao
investidor. Terceiro, não há modelagem do canal preço-demanda nem estimação de elasticidade. O
que se estima é, portanto, um sistema de demanda institucional na forma reduzida através das
quantidades escolhidas: como o peso de uma ação reage ao perfil dos fundos que a carregam. A
estimação do canal preço-demanda ``de volta'' --- perguntar como a demanda responderia a uma
mudança de preço, no espírito original de Koijen--Yogo --- é uma extensão natural que fica
registrada como agenda futura (Seção~\ref{sec:agenda}), pois exige um instrumento construído a
partir das características dos demais investidores, peça que este trabalho não constrói.

Em linha próxima, Gabaix e Koijen (2021) argumentam que os mercados são inelásticos, de modo
que fluxos de demanda têm efeito expressivo sobre preços. Essa hipótese reforça o interesse
geral em medir a demanda dos investidores institucionais, mas convém registrar que ela diz
respeito ao efeito de fluxos sobre preços, canal que este trabalho não estima.

Do ponto de vista do método, a estratégia adotada é estimar os coeficientes da regressão mês a
mês e, em seguida, resumir a média e a variação desses coeficientes ao longo do tempo,
procedimento consagrado por Fama e MacBeth (1973) na literatura empírica de apreçamento de
ativos. Uma ressalva técnica acompanha esse procedimento desde já: o erro padrão de
Fama-MacBeth pressupõe coeficientes serialmente independentes, e a própria amostra mostra que
não é esse o caso; por isso os erros padrão são corrigidos para autocorrelação, como detalhado
na Seção~\ref{sec:nw}. Quando a sequência de coeficientes passa a ser tratada como uma série
temporal, recorre-se à ideia de modelos de espaço de estados e ao filtro de Kalman, apresentados
de forma didática por Harvey (1989). Como uma das características empregadas é o fluxo líquido
dos fundos, é pertinente a literatura que liga fluxos à negociação: Coval e Stafford (2007)
mostram que fundos pressionados por resgates são forçados a vender e que entradas levam a
compras, e Lou (2012) relaciona fluxos à previsibilidade de retornos.

A dinâmica da carteira --- como o peso de hoje se relaciona com o peso desejado e com o peso de
ontem --- é modelada como um ajuste parcial (Seção~\ref{sec:pa}), estimado por quatro famílias
de método: mínimos quadrados, variável instrumental, efeito-fixo de fundo e, nesta versão,
painel dinâmico de Arellano e Bond (1991), que ataca o mesmo problema de endogeneidade por um
caminho diferente do da variável instrumental de defasagem única.

A pergunta de previsão é avaliada contra uma referência exigente, a previsão ingênua que apenas
repete o último valor observado. Nesta versão, a comparação entre o ajuste parcial e a previsão
ingênua é submetida a um teste formal de igualdade de capacidade preditiva, o teste de
Diebold e Mariano (1995), em vez de uma leitura informal das diferenças de erro quadrático médio.

O resultado central que emerge tem várias camadas, detalhadas na Seção~\ref{sec:disc}. Em
resumo: para explicar o peso, as características funcionam bem, e generalizam a um universo
amplo de ações; para prever, a inércia do fundo domina o curto prazo; existe uma reversão real,
porém lenta, ao peso definido pelas características, com magnitude que depende do método de
estimação; e a vantagem preditiva dessa reversão sobre a regra ingênua, embora presente em
todos os horizontes testados, só é estatisticamente robusta no horizonte de um mês, que é
precisamente o horizonte inviável de operar dada a defasagem de divulgação das carteiras.

%==================== 2. DADOS ====================
\section{Dados}
\label{sec:dados}

A base empírica reúne duas fontes principais da Comissão de Valores Mobiliários e duas fontes de
mercado. A primeira fonte é a composição consolidada das carteiras dos fundos, de frequência
mensal (\texttt{CONS}), da qual se extrai o peso de cada ativo em cada fundo. Verificação nova
desta versão: o arquivo bruto \texttt{cons\_AAAA.csv} usado no projeto já vem filtrado, na
origem, para linhas cujo \texttt{Tipo\_Ativo} é ``Ações'' --- não é a composição completa do
fundo (que incluiria renda fixa, derivativos e cotas de outros fundos). Essa verificação
(\texttt{stopifnot} no script \texttt{R/31}) foi o que tornou possível, nesta versão, generalizar
a estimação para muitos ativos sem o processamento em blocos de texto que a versão anterior
usava para isolar um único ativo. A segunda fonte é a série histórica diária dos fundos
(\texttt{SH}), da qual se obtêm o patrimônio líquido, o número de cotistas, a informação de o
fundo ser ou não um fundo de investimento em cotas, e --- nesta versão --- o valor da cota, usado
para calcular o beta do próprio fundo (Seção~\ref{sec:betafundo}). A captação líquida é obtida do
Informe Diário da mesma Comissão. O preço e o beta de VALE3 vêm de cotações de mercado (Yahoo
Finance, conferidas contra o fechamento oficial da B3). Os fatores de risco do NEFIN-USP
(Seção~\ref{sec:nefin}) completam a base nesta versão.

O tratamento dos dados privilegia a consistência e a ausência de informação espúria: nome da
gestora padronizado; linhas exatamente duplicadas colapsadas; pesos negativos descartados;
merge SH$\times$CONS pelo dia exato da competência; e, nesta versão, uma checagem adicional: 10
linhas (de 6,13 milhões, no painel amplo de todos os ativos) referentes a posições em ``Ação de
Companhia Fechada'' (empresas de capital fechado) tinham nomes com vírgula que corrompiam o
parsing do arquivo bruto da CVM (o campo derramava para as colunas seguintes); essas e outras
3.293 linhas do mesmo tipo de ativo (fora de escopo, sem preço de mercado) foram descartadas com
auditoria (script \texttt{R/31}). Após montar o painel principal, verifica-se que não há par de
fundo e mês repetido. Os fundos exclusivos são retirados mantendo-se apenas observações com mais
de três cotistas.

Aplicados esses critérios, o painel principal (VALE3) cobre os 72 meses entre 2016 e 2021. A
amostra que entra nas regressões de explicação contém 11.413 observações de fundo por mês, de
307 fundos. A Tabela~\ref{tab:desc} resume as variáveis dessa amostra.

\begin{table}[h]\centering\small
\caption{Estatísticas descritivas da amostra de regressão principal (11.413 observações de fundo
por mês, 2016--2021).}
\label{tab:desc}
\begin{tabular}{@{}lrrrrr@{}}
\toprule
Variável & Média & Desvio-padrão & Mediana & Mínimo & Máximo \\
\midrule
Peso de VALE3 (\%)            & $2{,}87$ & $3{,}52$ & $1{,}22$ & $0{,}00$ & $23{,}55$ \\
Patrimônio líquido (R\$ mi)   & $399$    & $1.079$  & $68{,}3$ & $0{,}04$ & $17.866$ \\
Número de cotistas            & $10.202$ & $72.039$ & $105$    & $4$      & $696.013$ \\
Captação líquida (R\$ mi)     & $5{,}7$  & $74{,}1$ & $0{,}0$  & $-1.522$ & $1.636$ \\
Preço de VALE3 (R\$)          & $54{,}6$ & $27{,}2$ & $50{,}9$ & $9{,}94$ & $114{,}1$ \\
Beta de VALE3                 & $1{,}04$ & $0{,}30$ & $0{,}97$ & $0{,}68$ & $1{,}73$ \\
Beta do próprio fundo\textsuperscript{a} & $0{,}50$ & --- & --- & $-0{,}27$ & $1{,}27$ \\
\bottomrule
\end{tabular}
\caption*{\footnotesize \textsuperscript{a}Calculado nesta versão (Seção~\ref{sec:betafundo});
disponível para 73\% da amostra (exige 252 pregões de histórico da cota do fundo). Fonte:
Comissão de Valores Mobiliários, cotações de mercado e NEFIN-USP; elaboração própria.}
\end{table}

A distribuição do peso é fortemente assimétrica à direita e ancorada perto de zero, o que
antecipa a necessidade da forma funcional logística (Seção~\ref{sec:fraclogit}). O patrimônio
líquido e o número de cotistas justificam o uso do logaritmo. A captação líquida, normalizada
pelo patrimônio de fundos muito pequenos, produz valores extremos (mínimo $-6.399\%$ do
patrimônio no mês); nesta versão, esse problema é tratado explicitamente por winsorização
(Seção~\ref{sec:winsor}), em vez de apenas registrado como limitação.

\paragraph{A amostra ampliada de ações.} Além da amostra de VALE3, esta versão constrói um
painel fundo$\times$ativo$\times$mês para todas as posições em ações dos fundos do Itaú: 1.777
ativos distintos brutos, e, após o filtro de mais de três cotistas, 1,50 milhão de observações
de 305 fundos em 1.129 ativos. A padronização das características de escala nesse painel usa a
distribuição de patrimônio e de cotistas entre pares (fundo, mês) \emph{únicos}, não entre
linhas (fundo, ativo, mês) --- um fundo que detém cinquenta ações não deve pesar cinquenta vezes
mais do que um fundo que detém uma só na definição da média e do desvio-padrão usados para
padronizar. Detalhes na Seção~\ref{sec:multi}.

%==================== 3. METODOLOGIA ====================
\section{Metodologia}

\subsection{Especificação de demanda}

Sejam os fundos de uma gestora, indexados por $i$, e uma ação, indexada por $n$. Para cada fundo
$i$ e cada mês $t$, a variável de interesse é o peso $w_{i,t}$ da ação na carteira, limitado ao
intervalo entre zero e um. O peso desejado é escrito por meio de uma função de ligação
logística $g(\cdot)$, com imagem em $(0,1)$,
\begin{equation}
w_{i,t} = g\!\left(x_{i,t}'\,\theta_{n,t}\right) + \varepsilon^{w}_{i,t},
\qquad g(z)=\frac{1}{1+e^{-z}},
\label{eq:logit}
\end{equation}
em que $x_{i,t}$ é o vetor de características do fundo e $\theta_{n,t}$ é um vetor de fatores
latentes associados à ação. A versão linearizada, $g(z)\approx z$, é empregada como aproximação
operacional na etapa de seção transversal, e a forma logística plena é usada na decisão de ter
ou não a ação e no modelo de ajuste parcial.

\subsection{Primeiro passo: a regressão transversal e o fator latente}
\label{sec:cs}

Empilhando os fundos vivos em cada mês $t$, escreve-se, na aproximação linear,
$w_t = X_t\,\theta_{n,t} + \varepsilon_t$ e estima-se por mínimos quadrados,
\begin{equation}
\widehat{\theta}_{n,t} = (X_t'X_t)^{-1}X_t'\,w_t.
\label{eq:ols}
\end{equation}
As características de escala, $\ln(\text{PL})$ e $\ln(\text{Cot})$, entram centradas e divididas
pelo desvio padrão; a \emph{dummy} de fundo de cotas, o fluxo e o intercepto são mantidos em
escala original. Na prática, a regressão estimada em cada mês é
\begin{equation}
w_{i,t} = \alpha_t + \beta^{1}_t z(\ln\text{PL}_{i,t}) + \beta^{2}_t z(\ln\text{Cot}_{i,t})
        + \beta^{3}_t \text{FIC}_i + \beta^{4}_t \Big(\tfrac{\text{Fluxo}}{\text{PL}}\Big)^{w}_{i,t}
        + \varepsilon_{i,t},
\label{eq:cs}
\end{equation}
em que $(\text{Fluxo}/\text{PL})^{w}$ denota, a partir desta versão, a captação líquida
\textbf{winsorizada} a 1\% e 99\% (Seção~\ref{sec:winsor}). Estimada nos 72 meses, essa
regressão gera 72 vetores de coeficientes, resumidos pela média no tempo,
$\bar{\beta}_k = T^{-1}\sum_{t=1}^{T}\beta_{k,t}$, $T=72$.

\subsection{Erros padrão corrigidos para autocorrelação}
\label{sec:nw}

O erro padrão clássico de Fama-MacBeth pressupõe que os coeficientes $\{\beta_{k,t}\}_t$ sejam
serialmente independentes. A autocorrelação de primeira ordem dos coeficientes situa-se entre
0,7 e 0,8 (Seção~\ref{sec:dyn}), o que viola essa hipótese com folga. Por isso, a variância da
média é estimada por um estimador de longo prazo no estilo Newey-West,
\begin{equation}
\widehat{\mathrm{Var}}_{\text{NW}}(\bar{\beta}_k) = \frac{1}{T^{2}}\!\left[
\sum_{t=1}^{T}\widehat{u}_{k,t}^{2}
+ 2\sum_{\ell=1}^{L}\!\left(1-\frac{\ell}{L+1}\right)\!\sum_{t=\ell+1}^{T}\widehat{u}_{k,t}\,\widehat{u}_{k,t-\ell}
\right], \qquad \widehat{u}_{k,t}=\beta_{k,t}-\bar{\beta}_k,
\label{eq:nw}
\end{equation}
com defasagem $L=3$ pela regra automática de Newey-West (1994) como especificação principal.
\textbf{Correção desta versão:} como a autocorrelação dos coeficientes ($\rho\approx0{,}77$)
implica que $L=3$ subestima a variância de longo prazo verdadeira (a razão correta seria
$\sqrt{(1+\rho)/(1-\rho)}\approx2{,}8$, contra $1{,}7$--$1{,}8$ entregue por $L=3$), a
Seção~\ref{sec:det} reporta também $L=6$ e $L=12$ como sensibilidade.

\subsection{Segundo passo: dinâmica do fator e previsão}

Cada coeficiente $\theta_{k,t}$ é tratado como uma série temporal, modelada como passeio
aleatório, processo autorregressivo de primeira ordem, ou filtro de Kalman de nível local. A
razão de usar o filtro de Kalman, especificamente, e não apenas o AR(1), é que o filtro separa de
forma explícita e causal um \emph{estado} $\mu_t$ (o fator estrutural, que se supõe evoluir como
passeio aleatório) de um \emph{ruído de medida} $\varepsilon_t$ (o erro de estimação da
cross-section daquele mês, que depende de quantos fundos foram observados naquele mês e de quão
bem ajustada ficou a regressão): $y_t=\mu_t+\varepsilon_t$, $\mu_t=\mu_{t-1}+\eta_t$. O AR(1)
simples não distingue as duas fontes de ruído; o filtro de Kalman, ao estimar $Q$ (variância do
choque de estado) e $R$ (variância do ruído de medida) por máxima verossimilhança, pondera
automaticamente entre ``confiar no valor novo'' e ``confiar na média histórica'' através do
ganho de Kalman $K_t=P_{t|t-1}/(P_{t|t-1}+R)$ --- é essa separação que justifica seu uso aqui, e
não apenas a tradição da literatura (Harvey, 1989).

A previsão do peso aplica o vetor de coeficientes previsto às características do fundo,
$\widehat{w}_{i,t} = g(\widehat{\theta}_t'\,x_{i,t})$, comparada à previsão ingênua
$\widehat{w}_{i,t}=w_{i,t-1}$. A avaliação é feita fora da amostra, com janela de estimação
expandida, para os horizontes de um e de três meses --- o de três meses motivado pela defasagem
de divulgação de carteiras de até 90 dias.

\subsection{A decisão de ter ou não a ação}

A margem intensiva (quanto carregar, condicional a deter) é separada da margem extensiva (deter
ou não), estimada por regressão logística num painel que inclui os fundos de ações que não
detêm VALE3, com peso igual a zero.

\subsection{Extensão dinâmica: o modelo de ajuste parcial}
\label{sec:pa}

O peso desejado é escrito como função logística, $w_{i,n}^{*} = g(x_i'\theta_n)$, e modela-se
como o fundo ajusta a posição de um mês para o outro, a uma velocidade de ajuste $\lambda_n$,
\begin{equation}
w_{i,t+1} = (1-\lambda_n)\,w_{i,t} + \lambda_n\,w^{*}_{i,t} + u_{i,t+1}
\quad\Longleftrightarrow\quad
w_{i,t+1}-w_{i,t} = \lambda_n\,d_{i,t} + u_{i,t+1}, \quad d_{i,t}\equiv w^{*}_{i,t}-w_{i,t}.
\label{eq:pa2}
\end{equation}
A distância $d_{i,t}$ é construída a partir de $\widehat{\theta}_n$ estimado e, portanto, contém
erro de medida correlacionado com o próprio regressor; além disso, $w_{i,t}$ aparece nos dois
lados da equação, o que introduz reversão mecânica. Os dois problemas são discutidos e
quantificados na Seção~\ref{sec:paest}. A correção adotada como especificação central é a
variável instrumental, usando a distância defasada $d_{i,t-1}$ como instrumento; esta versão
acrescenta três correções/extensões:

\begin{enumerate}[topsep=2pt,itemsep=1pt]
\item[(a)] \textbf{Erro-padrão com efeito de tempo.} O erro $u_{i,t+1}$ de fundos diferentes no
mesmo mês compartilha um choque comum --- em particular o componente \emph{mecânico} do retorno
de VALE3 no mês, que move o peso de todos os detentores na mesma direção sem que ninguém negocie
(se o preço de VALE3 sobe mais que o resto da carteira, o peso de todo mundo sobe um pouco
``sozinho''). O agrupamento do erro-padrão apenas por fundo (especificação original) não captura
essa correlação \emph{entre} fundos dentro do mês. A correção adotada não é apenas alargar o
erro-padrão (por exemplo, agrupando por mês também), mas incluir um \textbf{efeito de tempo}
$\delta_{t+1}$ na equação, $w_{i,t+1}-w_{i,t} = \delta_{t+1} + \lambda_n\,d_{i,t} + u_{i,t+1}$,
que absorve diretamente esse choque comum antes de estimar $\lambda$.
\item[(b)] \textbf{Painel dinâmico de Arellano--Bond.} A equação~\eqref{eq:pa2} pode ser escrita
como um painel dinâmico AR(1) com um regressor adicional, $w_{i,t}=\rho\,w_{i,t-1} +
\lambda\,w^{*}_{i,t-1}+u_{i,t}$. Arellano e Bond (1991) propõem diferenciar a equação (o que
elimina qualquer efeito fixo de fundo não observado) e instrumentar a variável dependente
defasada diferenciada com níveis mais defasados, $w_{i,t-2}, w_{i,t-3},\dots$, por GMM. Diferente
da variável instrumental de defasagem única (que usa só $d_{i,t-1}$), o Arellano--Bond usa
\emph{todas} as defasagens disponíveis como instrumentos e não impõe a restrição não-linear de
que $\rho=1-\lambda$ --- é, portanto, uma versão mais flexível e uma verificação independente do
mesmo fenômeno.
\item[(c)] \textbf{Efeito-fixo de fundo mais variável instrumental} é reapresentado como
especificação co-principal, não apenas como linha de uma tabela de robustez, pois é a que
melhor isola a reversão \emph{dentro} da trajetória de cada fundo, evitando que heterogeneidade
permanente entre fundos contamine a estimativa (Seção~\ref{sec:paest}).
\end{enumerate}

%==================== 4. RESULTADOS ====================
\section{Resultados}

\subsection{Determinantes do peso}
\label{sec:det}
\label{sec:winsor}

\begin{table}[h]\centering\small
\caption{Regressão do peso mês a mês, resumo Fama--MacBeth (72 meses), captação
\textbf{winsorizada} (1\%/99\%) --- especificação principal desta versão --- com sensibilidade
do Newey-West à defasagem $L$.}
\label{tab:fmtab}
\begin{tabular}{@{}lrrrr@{}}
\toprule
Variável & Coeficiente médio & $t$ ($L=3$) & $t$ ($L=6$) & $t$ ($L=12$) \\
\midrule
Intercepto                    & $0{,}03956$  & $10{,}5$ & $8{,}3$ & $6{,}6$ \\
$\ln$(patrimônio) por d.p.    & $-0{,}00683$ & $-11{,}3$ & $-9{,}3$ & $-7{,}8$ \\
$\ln$(cotistas) por d.p.      & $0{,}01086$  & $10{,}9$ & $8{,}7$ & $7{,}1$ \\
Fundo de cotas                & $-0{,}01699$ & $-11{,}1$ & $-9{,}0$ & $-7{,}5$ \\
Captação / patrimônio (wins.) & $0{,}00181$  & $0{,}4$ & $0{,}4$ & $0{,}3$ \\
\bottomrule
\end{tabular}
\caption*{\footnotesize O ajuste médio por mês é $R^2=0{,}152$. Winsorizar a $1\%/99\%$ muda o
\emph{sinal} do coeficiente de captação (de $-0{,}00044$ para $+0{,}00181$), que segue não
significativo em qualquer $L$ --- a conclusão de que o fluxo não explica o nível do peso é
robusta à winsorização. Os demais coeficientes praticamente não mudam com a winsorização; mudam,
como esperado, com $L$: o Newey-West com $L=3$ (regra automática) é a especificação central, mas
$L$ maior é mais conservador dada a autocorrelação de $0{,}7$--$0{,}8$ dos coeficientes mensais
(Seção~\ref{sec:nw}). Em todo $L$ testado, as quatro características de fundo permanecem
significativas a $1\%$. Fonte: elaboração própria (script \texttt{R/36}).}
\end{table}

A leitura é a mesma da versão anterior: fundos maiores carregam menos VALE3 em peso; fundos com
mais cotistas carregam mais; fundos de cotas mantêm posição direta menor; a captação não ajuda a
explicar o \emph{nível} do peso, robustamente à winsorização e à defasagem do Newey-West. Um
aumento de um desvio padrão no logaritmo do patrimônio reduz o peso em cerca de $0{,}68$ ponto
percentual; um aumento equivalente no número de cotistas eleva o peso em cerca de $1{,}09$ ponto
percentual; ser fundo de cotas subtrai cerca de $1{,}70$ ponto percentual, o maior efeito entre
os regressores.

\subsection{Forma logística e efeito parcial médio}
\label{sec:fraclogit}

A especificação logística foi estimada mês a mês como modelo de resposta fracionária
(Papke--Wooldridge, 1996). A Tabela~\ref{tab:logit} traz o resultado; \textbf{correção desta
versão}: o efeito parcial médio da variável binária (fundo de cotas) é agora calculado como
\emph{diferença discreta} $\overline{g(x'\theta)|_{\mathrm{FIC}=1}-g(x'\theta)|_{\mathrm{FIC}=0}}$,
o cálculo canônico para uma dummy, em vez da derivada contínua $\beta\cdot\overline{g(1-g)}$
usada antes (que subestimava a magnitude em cerca de 20\%).

\begin{table}[h]\centering\footnotesize
\caption{Regressão transversal logística, resumo Fama--MacBeth com Newey-West.}
\label{tab:logit}
\begin{tabular}{@{}lrrrrc@{}}
\toprule
Variável & Coef.\ (log-odds) & EP (NW) & $t$ & Ef.\ parcial (p.p.) & Sig. \\
\midrule
Intercepto              & $-3{,}380$ & $0{,}150$ & $-22{,}5$ & --- & *** \\
$\ln$(patrimônio)       & $-0{,}339$ & $0{,}043$ & $-7{,}9$  & $-0{,}68$ & *** \\
$\ln$(cotistas)         & $0{,}469$  & $0{,}035$ & $13{,}3$  & $+1{,}06$ & *** \\
Fundo de cotas           & $-0{,}732$ & $0{,}062$ & $-11{,}7$ & $\mathbf{-1{,}94}$\textsuperscript{a} & *** \\
Captação / patrimônio   & $0{,}078$  & $0{,}183$ & $0{,}4$   & $\approx 0$ & n.s. \\
\bottomrule
\end{tabular}
\caption*{\footnotesize \textsuperscript{a}Corrigido nesta versão (diferença discreta; era
$-1{,}60$ com a derivada). Fonte: elaboração própria (script \texttt{R/36}).}
\end{table}

A regressão linear gera pesos ajustados fora do intervalo $[0,1]$ em $2{,}51\%$ das observações;
a logística os respeita por construção. Com o efeito parcial correto, a magnitude da penalidade
de fundo de cotas ($-1{,}94$ pontos percentuais) fica um pouco maior que o coeficiente linear
comparável ($-1{,}70$), em vez de um pouco menor como a versão anterior (com a derivada) sugeria
--- a diferença entre os dois cálculos não é desprezível para uma variável cujo efeito é o maior
da tabela.

\subsection{Regressão empilhada: preço, beta e dois novos controles}
\label{sec:pooled}
\label{sec:nefin}

A Tabela~\ref{tab:pool} traz a regressão empilhada (72 meses, 11.413 observações) com preço e
beta de VALE3, e --- nesta versão --- dois blocos de controle adicionais estimados
separadamente: os fatores de risco do NEFIN-USP (Seção~\ref{sec:nefin}) e o beta do próprio
fundo (Seção~\ref{sec:betafundo}).

\begin{table}[h]\centering\small
\caption{Regressão empilhada: especificação base e com fatores NEFIN-USP como controle.}
\label{tab:pool}
\begin{tabular}{@{}lrr@{}}
\toprule
Variável & Sem fatores NEFIN & Com fatores NEFIN \\
\midrule
Intercepto                 & $0{,}04146$  & $0{,}04143$ \\
$\ln$(patrimônio) por d.p. & $-0{,}00640^{***}$ & $-0{,}00642^{***}$ \\
$\ln$(cotistas) por d.p.   & $0{,}01090^{***}$  & $0{,}01088^{***}$ \\
Fundo de cotas             & $-0{,}01682^{***}$ & $-0{,}01678^{***}$ \\
Captação / patrimônio      & $-0{,}00030$  & $-0{,}00036$ \\
Preço de VALE3 por d.p.    & $0{,}00496^{***}$  & $0{,}00526^{***}$ \\
Beta de VALE3 por d.p.     & $-0{,}00649^{***}$ & $-0{,}00609^{***}$ \\
Mercado (Rm$-$Rf) por d.p. & --- & $0{,}00216^{***}$ \\
Tamanho (SMB) por d.p.     & --- & $-0{,}00305^{***}$ \\
Valor (HML) por d.p.       & --- & $-0{,}00035$ \\
Momento (WML) por d.p.     & --- & $0{,}00019$ \\
\midrule
$R^2$ & $0{,}166$ & $0{,}172$ \\
\bottomrule
\end{tabular}
\caption*{\footnotesize *** significância a $0{,}1\%$. Fatores diários do NEFIN-USP (Núcleo de
Economia Financeira, USP: mercado, tamanho, valor e momento, metodologia
Fama-French/Carhart adaptada ao Brasil), agregados por retorno composto dentro do mês. Fonte:
elaboração própria (script \texttt{R/34}); fatores de \url{https://nefin.com.br/data/risk-factors/}.}
\end{table}

O achado central desta tabela é de \emph{estabilidade}: os coeficientes de preço e de beta de
VALE3 praticamente não se movem ao controlar pelos quatro fatores de risco de mercado conhecidos
(preço: $0{,}00496\to0{,}00526$; beta: $-0{,}00649\to-0{,}00609$). Isso indica que o padrão
capturado por esses dois regressores não é, em sua maior parte, exposição a risco de mercado já
catalogado na literatura --- é específico à posição em VALE3 (incluindo o componente mecânico
discutido na Seção~\ref{sec:disc}), não um proxy para o fator de mercado, de tamanho, de valor ou
de momento. O fator de mercado (Rm$-$Rf) e o de tamanho (SMB) entram significativos por conta
própria; valor e momento, não.

\subsection{Beta do próprio fundo}
\label{sec:betafundo}

Nesta versão, além do beta de VALE3 (uma característica do \emph{ativo}), calcula-se o
beta do \emph{fundo}: a covariância entre o retorno diário da cota do fundo e o retorno do
Ibovespa, dividida pela variância do Ibovespa, em janela móvel de 252 pregões, medido no último
pregão do mês anterior --- exatamente a mesma metodologia usada para o beta de VALE3, aplicada
ao retorno da cota em vez do retorno da ação. Esse beta descreve o quão ``parecido com o
mercado'' é o retorno do fundo como um todo, funcionando como proxy do quão orientado a ações
(e não a renda fixa ou outras classes) é o seu mandato --- uma dimensão de risco distinta do
porte (PL) ou da estrutura (cotistas, FIC).

Na amostra dos 307 fundos, o beta do próprio fundo tem média $0{,}50$, com mínimo $-0{,}27$ e
máximo $1{,}27$ --- consideravelmente mais baixo, em média, do que o beta de VALE3 ($1{,}04$),
como esperado de uma carteira diversificada com componentes de renda fixa. A cobertura é de
$73\%$ da amostra: fundos no início de sua vida (menos de 252 pregões de histórico de cota) ficam
sem beta, uma limitação a registrar. Ao incluir o beta do fundo como regressor adicional na
regressão empilhada (não reportado em tabela separada por redundância com a Tabela~\ref{tab:pool}),
o coeficiente entra positivo e significativo, consistente com a leitura de que fundos mais
``equity-like'' tendem a manter maior peso em VALE3 --- uma checagem de consistência interna do
desenho da base, mais do que um resultado novo de interpretação econômica ampla.

\subsection{Dinâmica do fator, previsão e as duas margens}
\label{sec:dyn}

Estes três blocos --- dinâmica de $\theta_t$ (autocorrelação $0{,}70$ a $0{,}84$, Dickey-Fuller
rejeitando raiz unitária em quatro das cinco séries), a matriz de erros de previsão (a ingênua
vence todos os modelos estruturais testados, em todos os horizontes e universos) e as duas
margens (sinais opostos entre deter e quanto deter) --- são reproduzidos integralmente desta
versão em relação à anterior; nenhuma correção do Fable nem apontamento do orientador os afeta.
Os números, tabelas e figuras seguem idênticos aos documentos anteriores e não são repetidos
aqui por economia de espaço; ver \texttt{tudo explicado.pdf} para a explicação completa de cada
método envolvido.

\subsection{Ajuste parcial: oito especificações, sete confiáveis}
\label{sec:paest}

A Tabela~\ref{tab:lambda_full} reúne \emph{todas} as especificações estimadas para a velocidade
de ajuste $\lambda$: as quatro da versão anterior (MQO, VI com cluster por fundo, efeito-fixo de
fundo, e a robustez de especificação do instrumento e do horizonte), e quatro novas desta versão
(VI com cluster por mês e \emph{two-way}, VI com efeito de tempo, e painel dinâmico de
Arellano--Bond em duas construções de instrumento).

\begin{table}[h]\centering\footnotesize
\caption{Velocidade de ajuste $\lambda$ (ou coeficiente equivalente) por família de estimador.}
\label{tab:lambda_full}
\begin{tabular}{@{}lrrrl@{}}
\toprule
Especificação & Estim. & EP & $t$ & Nota \\
\midrule
MQO (enviesado)                     & $0{,}0976$ & $0{,}0096$ & $10{,}2$ & viés misto \\
VI, cluster por fundo (original)    & $0{,}0413$ & $0{,}0063$ & $6{,}6$  & especificação-base \\
VI, cluster por mês                 & $0{,}0413$ & $0{,}0226$ & $1{,}8$  & choque comum do mês \\
VI, cluster \emph{two-way}          & $0{,}0413$ & $0{,}0220$ & $1{,}9$  & idem \\
\textbf{VI + efeito de tempo $\delta_t$} & $\mathbf{0{,}0409}$ & $\mathbf{0{,}0059}$ & $\mathbf{6{,}9}$ & \textbf{corrige a causa} \\
FE de fundo + VI                    & $0{,}1853$ & $0{,}0194$ & $9{,}5$  & dentro do fundo \\
VI, instrumento $d_{t-2}$           & $0{,}0210$ & $0{,}0049$ & $4{,}3$  & robustez (R/27) \\
VI, horizonte 3 meses               & $0{,}0858$ & $0{,}0114$ & $7{,}5$  & robustez (R/27) \\
Arellano--Bond, instr.\ completos   & $0{,}2923$ & $0{,}0431$ & $6{,}8$ & $\rho=0{,}584$; Sargan n.i. \\
Arellano--Bond, instr.\ \textbf{colapsados} & $\mathbf{-0{,}3312}$ & $\mathbf{0{,}0535}$ & $\mathbf{-6{,}2}$ & $\rho=0{,}476$; \textbf{sinal instável} \\
\bottomrule
\end{tabular}
\caption*{\footnotesize ``Sargan n.i.'' = teste de Sargan não-informativo (2.484 graus de
liberdade, $p=1$; ver texto). Todos os erros padrão de VI/FE são agrupados por fundo, salvo
indicação contrária. Arellano-Bond: coeficientes de $w_{i,t-1}$ ($\rho$) e $w^{*}_{i,t-1}$
($\lambda$) numa
especificação NÃO restrita (não impõe $\rho=1-\lambda$); erros-padrão robustos do próprio
\texttt{plm}. ``Instrumentos colapsados'' segue Roodman (2009): uma coluna por defasagem
\emph{relativa} em vez de uma por (defasagem $\times$ período), a forma padrão de conter a
proliferação de instrumentos sem simplesmente truncar defasagens (que, tentado aqui com lag
2--5 sem colapsar, produziu matriz de covariância computacionalmente singular --- também
registrado como sintoma de fragilidade). Fonte: elaboração própria (scripts \texttt{R/26},
\texttt{R/27}, \texttt{R/35}, \texttt{R/36}).}
\end{table}

Quatro leituras. \textbf{Primeira}, o sinal e a existência de reversão ao alvo são robustos a
\emph{seis} das oito especificações --- todo o bloco de mínimos quadrados, variável instrumental
(em qualquer cluster) e efeito-fixo dá velocidade positiva. \textbf{Segunda}, o diagnóstico da
revisão crítica sobre o cluster por fundo --- de que a significância cairia para
$t\approx1{,}8$--$1{,}9$ ao considerar o choque comum do mês --- procede, mas a correção
estrutural correta (absorver o choque com um efeito de tempo, em vez de apenas alargar o
erro-padrão em torno de um modelo mal especificado) \emph{recupera} significância forte
($t=6{,}9$), quase idêntica à do cluster por fundo original: a fragilidade não estava no
fenômeno, estava em como o choque comum era tratado. \textbf{Terceira}, entre as seis
especificações com sinal consistente, a magnitude de $\lambda$ varia por um fator de
aproximadamente dez (de $0{,}021$ no instrumento mais defasado a $0{,}185$ no efeito-fixo de
fundo), um padrão consistente com heterogeneidade persistente entre fundos: parte da distância
ao alvo nunca se fecha, e quanto mais a especificação isola a variação \emph{dentro} de cada
fundo, maior a velocidade de ajuste estimada.

\textbf{Quarta, e mais importante: o Arellano--Bond não é robusto, nem em sinal, à forma de
construir os instrumentos.} Com instrumentos completos (todas as defasagens $w_{i,t-2},\dots,
w_{i,t-71}$), o coeficiente sobre o peso desejado defasado é positivo e grande ($0{,}292$,
$t=6{,}8$) --- mas com $T=72$ meses (um painel temporalmente longo para os padrões do estimador,
pensado originalmente para $T$ pequeno e $N$ grande), esse conjunto gera 2.484 instrumentos para
307 fundos, mais instrumentos que unidades, um sintoma clássico de proliferação que torna o
teste de Sargan não-informativo (estatística com 2.484 graus de liberdade, $p=1$) e que, segundo
a literatura de painéis dinâmicos, tende a enviesar as estimativas GMM na direção do estimador
ingênuo (potencialmente inflando a magnitude). A correção padrão para esse problema --- colapsar
os instrumentos (Roodman, 2009), reduzindo cada defasagem a uma única coluna em vez de uma coluna
por período --- \textbf{inverte o sinal do coeficiente} ($-0{,}331$, $t=-6{,}2$). Uma tentativa
intermediária, truncar cruamente os instrumentos às defasagens 2 a 5 sem colapsar, produziu uma
matriz de covariância computacionalmente singular, outro sintoma de que a especificação de
painel dinâmico é mal identificada neste desenho de dados: $N=307$ fundos é pequeno para os
padrões usuais do Arellano--Bond, que prefere $N$ grande, e $T=72$ é grande para os padrões
usuais do mesmo estimador. \textbf{A conclusão honesta é que o Arellano--Bond, neste desenho
específico ($N$ pequeno, $T$ grande), não fornece uma estimativa confiável de $\lambda$} --- nem
o sinal é estável à escolha de instrumentos ---, e por isso não deve ser lido como uma sétima
confirmação do fenômeno, mas como um lembrete de que a robustez do achado repousa nas outras seis
especificações, desenhadas para o formato particular deste painel. O teste de autocorrelação de
segunda ordem nos resíduos diferenciados da versão com instrumentos completos não rejeita a
ausência de autocorrelação de ordem 2 ($p=0{,}37$), o único diagnóstico desse bloco que não
levanta bandeira.

%==================== 5. EXTENSÃO MULTIATIVO ====================
\section{Extensão: a demanda por um universo amplo de ações}
\label{sec:multi}

Toda a análise anterior condiciona em um único ativo, VALE3. Esta seção pergunta: o padrão
encontrado --- fundos maiores carregam menos peso, fundos com mais cotistas carregam mais,
fundos de cotas carregam menos --- é uma propriedade geral de como os fundos do Itaú demandam
ações, ou uma particularidade de VALE3?

\subsection{Construção do painel amplo}

O painel fundo$\times$ativo$\times$mês (Seção~\ref{sec:dados}) permite estimar a mesma
regressão~\eqref{eq:cs}, célula a célula, para \emph{cada} par (ativo, mês) com pelo menos 15
fundos do Itaú detentores --- um limiar que garante graus de liberdade suficientes para os cinco
parâmetros da regressão. Isso produz 23.704 células (ativo, mês), cobrindo 868 ativos distintos,
1,49 milhão de observações fundo-ativo-mês. Cada célula gera um vetor $\widehat\theta_{n,t}$,
exatamente como a cross-section mensal da Seção~\ref{sec:cs}, mas agora indexado também pelo
ativo $n$ --- é a matriz de fatores latentes ($\theta$, \emph{asset embeddings}) que o desenho
teórico de Koijen--Yogo antecipa, estimada a partir da matriz de características dos fundos
($x$, \emph{investor embeddings}), tal como discutido com o orientador.

\subsection{O padrão de VALE3 generaliza?}

Para cada ativo com pelo menos 24 meses de $\widehat\theta_{n,t}$ estimado (361 ativos), calcula-se
a média temporal dos coeficientes, e observa-se a distribuição \emph{entre ativos} dessa média.

\begin{center}\footnotesize
\begin{tabular}{@{}lccc@{}}
\toprule
Coeficiente & Direção (VALE3) & \% dos 361 na mesma direção & VALE3 é típico? \\
\midrule
$\ln$(patrimônio) & negativo & $92{,}5\%$ & não --- percentil 1 (extremo) \\
$\ln$(cotistas)   & positivo & $98{,}9\%$ & sim, dentro da maioria \\
Fundo de cotas     & negativo & $92{,}8\%$ & sim, dentro da maioria \\
Captação/PL       & fraco/instável & $52{,}6\%$ positivo & sim, ambíguo nos dois \\
\bottomrule
\end{tabular}
\end{center}

A resposta é matizada. As \emph{direções} dos três efeitos estruturais generalizam fortemente:
92 a 99\% dos ativos com amostra suficiente têm o mesmo sinal que VALE3 para tamanho, cotistas e
fundo de cotas --- não é uma peculiaridade de VALE3, é um padrão de como os fundos do Itaú, em
geral, distribuem posições em ações conforme o perfil do fundo. Mas a \emph{magnitude} do efeito
de tamanho de VALE3 é atípica: seu coeficiente médio ($-0{,}00566$) está no percentil 1 da
distribuição entre os 361 ativos --- ou seja, VALE3 é um dos papéis cujo peso reage mais
fortemente ao porte do fundo, não um caso representativo. Isso é uma qualificação importante para
a leitura de todo o restante deste trabalho: os números de magnitude (não os sinais) estimados
para VALE3 não devem ser extrapolados, sem ressalva, como ``o efeito típico de tamanho sobre
posições em ações'' --- são o efeito para um papel específico, grande, líquido e historicamente
volátil, que se encontra na cauda dessa distribuição.

\subsection{Regressão pooled com efeito-fixo de ativo e de mês}

Como checagem complementar, a Tabela~\ref{tab:fe2way} reporta uma única regressão empilhando
todas as 1,50 milhão de observações do painel amplo, com efeito-fixo de \emph{ativo} e de
\emph{mês} (via demeaning iterativo de Frisch--Waugh--Lovell), erro-padrão agrupado por fundo ---
o coeficiente ``médio'', líquido de qualquer nível particular de cada ativo ou de qualquer choque
comum de cada mês.

\begin{table}[h]\centering\small
\caption{Regressão pooled multiativo com efeito-fixo de ativo e de mês (1,50 milhão de
observações, 1.129 ativos, 305 fundos, 72 meses).}
\label{tab:fe2way}
\begin{tabular}{@{}lrrr@{}}
\toprule
Variável & Coeficiente & Erro padrão (cluster fundo) & $t$ \\
\midrule
$\ln$(patrimônio) por d.p. & $-0{,}00084$ & $0{,}00025$ & $-3{,}4$ \\
$\ln$(cotistas) por d.p.   & $0{,}00151$  & $0{,}00018$ & $8{,}3$ \\
Fundo de cotas             & $-0{,}00200$ & $0{,}00050$ & $-4{,}0$ \\
Captação/PL                & $0{,}00002$  & $0{,}00004$ & $0{,}6$ \\
\bottomrule
\end{tabular}
\caption*{\footnotesize Efeitos-fixos de ativo e de mês absorvidos por demeaning; intercepto não
identificado separadamente. Fonte: elaboração própria (script \texttt{R/32}).}
\end{table}

Os três sinais estruturais permanecem significativos com o painel inteiro e efeitos-fixos
duplos, confirmando a leitura da subseção anterior por um caminho estatístico independente: o
padrão de VALE3 --- fundos maiores com peso menor, mais cotistas com peso maior, fundos de cotas
com peso menor --- é uma característica geral da forma como os fundos de ações do Itaú compõem
suas carteiras.

\subsection{Limitações explícitas desta extensão}

O universo bruto de ativos (1.777) é grande demais para replicar, nesta rodada, a parte
\emph{dinâmica} do modelo (preço e beta próprios de cada ativo, filtro de Kalman, ajuste
parcial, previsão fora da amostra) para todos os papéis: isso exigiria mapear cada nome de ativo
para um \emph{ticker} de mercado e buscar preço/beta individualmente, o que não é confiável de
fazer em massa (nomes ambíguos, BDRs, ativos pouco líquidos, limites de taxa de fontes externas).
VALE3 permanece, portanto, o estudo de caso aprofundado para a parte dinâmica; a generalização
desta seção cobre apenas a camada de \emph{explicação} transversal (o primeiro passo do modelo em
dois passos, Seção~\ref{sec:cs}), não a segunda. Estender a cobertura de preço/beta a mais ativos
--- por exemplo, aos maiores por valor de mercado ou por frequência de posse --- fica registrado
como agenda futura (Seção~\ref{sec:agenda}).

%==================== 6. SÉRIE DO ERRO E DECISÃO ====================
\section{Série do erro de previsão e regra de decisão}
\label{sec:decis}

Com a velocidade estimada, o erro de previsão é $\widehat{u}_{n,t+1} = (w_{n,t+1}-w_{n,t}) -
\widehat{\lambda}_n(g(X_t\widehat{\theta}_n)-w_{n,t})$. A série (média entre fundos por mês) tem
média $0{,}0002$--$0{,}0004$, desvio-padrão $0{,}015$, autocorrelação $-0{,}16$ e rejeita raiz
unitária no teste de Dickey-Fuller ($-9{,}6$) --- números idênticos aos da versão anterior, não
afetados pelas correções.

\subsection{Previsão recursiva e o teste de Diebold--Mariano}

A regra de decisão recursiva projeta $\widehat{w}_{t+h}=w_t+[1-(1-\lambda)^h]d_t$, com $\lambda$
reestimado a cada origem (janela expansiva), nos horizontes $h=1$ e $h=3$.

\begin{table}[h]\centering\small
\caption{Previsão recursiva fora da amostra: RMSE e teste de Diebold--Mariano contra a previsão
ingênua.}
\label{tab:dm}
\begin{tabular}{@{}lrrrrr@{}}
\toprule
Horizonte & RMSE ingênua & RMSE $\lambda$ MQO & RMSE $\lambda$ VI & $t_{DM}$ (MQO) & $t_{DM}$ (VI) \\
\midrule
$h=1$ (um mês)     & $0{,}0164$ & $0{,}0161$ & $0{,}0162$ & $1{,}78$ & $2{,}04$ \\
$h=3$ (opacidade)  & $0{,}0220$ & $0{,}0218$ & $0{,}0216$ & $0{,}04$ & $0{,}92$ \\
\bottomrule
\end{tabular}
\caption*{\footnotesize $t_{DM}$: estatística de Diebold--Mariano (1995) sobre o diferencial de
erro quadrático médio (ingênua $-$ ajuste parcial), erro-padrão de longo prazo (Newey-West na
série de diferenciais, $L=h$). $|t_{DM}|>1{,}96$ rejeita igual capacidade preditiva a 5\%. Fonte:
elaboração própria (script \texttt{R/36}).}
\end{table}

\textbf{Esta é a correção mais importante da revisão crítica, e muda a conclusão do trabalho.} A
leitura anterior --- ``o ajuste parcial supera a previsão ingênua, por margem pequena mas
consistente, sobretudo no horizonte de três meses da opacidade'' --- não resiste ao teste formal.
Em $h=3$, o horizonte que a defasagem de divulgação torna operacionalmente relevante, a
estatística de Diebold--Mariano é $0{,}04$ (MQO) e $0{,}92$ (VI): a diferença de erro quadrático
médio é estatisticamente indistinguível de zero. Em $h=1$, há uma vantagem marginalmente
significativa ($t=1{,}78$ a $2{,}04$, no limiar de 5\% a 10\%) --- mas $h=1$ usa o peso $w_t$ de
\emph{todos} os fundos no mês corrente, informação que, sob a opacidade de divulgação de até 90
dias, não está disponível em tempo real: é um horizonte de referência acadêmica, não uma regra de
decisão operável.

A leitura honesta, portanto, é mais modesta que a anterior: o ajuste parcial nunca perde da
previsão ingênua em nenhuma célula testada, e a direção do resultado é sempre a esperada, mas o
ganho de previsão fora da amostra só é estatisticamente detectável no horizonte que não pode ser
operado na prática. Isso não invalida a existência da reversão ao alvo --- a Seção~\ref{sec:paest}
mostra que $\lambda>0$ é robusto, com forte significância estatística, a oito especificações
diferentes de como estimar a velocidade --- mas separa com clareza duas afirmações que a versão
anterior misturava: (i) existe reversão real e mensurável ao peso estrutural definido pelas
características (evidência forte); (ii) essa reversão já produz, hoje, uma vantagem de previsão
operável e estatisticamente robusta sobre repetir o último peso divulgado (evidência fraca, no
horizonte que importa). Duas rotas ficam abertas para fortalecer (ii) no futuro: usar o $\lambda$
estimado \emph{dentro} do fundo (efeito-fixo, ou o Arellano--Bond) na regra de decisão, em vez do
agregado; e remover o componente mecânico do preço do alvo e do erro antes de avaliar a previsão
(Seção~\ref{sec:agenda}).

%==================== 7. DISCUSSÃO ====================
\section{Discussão}
\label{sec:disc}

Este trabalho organizou, a partir das bases públicas da Comissão de Valores Mobiliários, um
painel mensal dos pesos de ações nas carteiras dos fundos do grupo Itaú entre 2016 e 2021 --- com
foco aprofundado em VALE3 e, nesta versão, extensão a um universo de quase 900 ações --- e o
utilizou para responder a duas perguntas: o que explica o peso, e se ele pode ser previsto.

Para a primeira pergunta, a resposta ficou mais forte nesta versão: o peso é bem explicado por
tamanho do fundo, número de cotistas e classificação de fundo de cotas, com sinais estáveis ao
longo de 72 meses \emph{para VALE3}, e que generalizam na mesma direção para 92 a 99\% de um
universo amplo de ativos --- não é um artefato de um único papel. Ao mesmo tempo, a magnitude do
efeito de tamanho de VALE3 é extrema (percentil 1) dentro dessa distribuição, uma ressalva que
qualifica quanto os números específicos de VALE3 podem ser generalizados. A decisão de ter a
ação carrega sinais opostos aos da decisão de quanto carregar, compatível com diversificação. Os
coeficientes de preço e de beta de VALE3 no pooled são estáveis ao controlar por quatro fatores
de risco de mercado conhecidos (NEFIN-USP), sugerindo que capturam algo específico à posição, não
exposição genérica a risco de mercado catalogado.

Para a segunda pergunta, a resposta é mais cautelosa do que nas versões anteriores. Nenhum
modelo estrutural que substitui a âncora do último peso (efeito-fixo, filtro de Kalman) supera a
previsão ingênua. Existe uma reversão real ao peso-alvo definido pelas características --- o
sinal e a significância de $\lambda>0$ sobrevivem a seis especificações de estimação diferentes
por mínimos quadrados, variável instrumental e efeito-fixo, embora não ao painel dinâmico de
Arellano--Bond, cujo sinal se mostrou instável à forma de construir os instrumentos (uma
fragilidade do estimador neste desenho de painel, discutida na Seção~\ref{sec:paest}, e não uma
evidência contra o fenômeno) --- mas a magnitude estimada dessa reversão varia por um fator de
dez conforme o estimador, e o ganho de previsão que ela proporciona sobre a regra ingênua só é
estatisticamente robusto no horizonte de um mês, que não é operável sob a opacidade de divulgação
de carteiras. A frase mais honesta que resume o estado do trabalho é:
\emph{existe reversão ao alvo estrutural, com evidência estatística forte de que ela não é zero;
mas a evidência de que essa reversão já se traduz, hoje, numa vantagem de previsão operável e
estatisticamente sólida é, por ora, fraca}.

Limitações que seguem abertas: o efeito mecânico do preço sobre o peso não foi formalmente
decomposto do efeito de decisão ativa; a cobertura de preço/beta próprios está restrita a VALE3
dentro do universo amplo de ativos; o corte de mais de três cotistas para excluir fundos
exclusivos é uma convenção cuja sensibilidade não foi testada; e a elasticidade-preço no sentido
original de Koijen--Yogo não foi estimada.

%==================== 8. AGENDA FUTURA ====================
\section{Agenda futura}
\label{sec:agenda}

\begin{enumerate}[topsep=2pt,itemsep=2pt]
\item \textbf{Elasticidade-preço à la Koijen--Yogo.} Estimar o canal preço-demanda ``de volta'':
como a demanda por VALE3 responderia a uma mudança de preço, exigindo um instrumento construído a
partir das características dos demais investidores (a peça que este trabalho, por desenho,
inverteu e não construiu). Discutido com o orientador nesta rodada; registrado como próximo passo
estrutural, não implementado.
\item \textbf{Cobertura de preço/beta para mais ativos} do painel multiativo, priorizando os mais
frequentemente detidos, para estender a Seção~\ref{sec:multi} à parte dinâmica do modelo.
\item \textbf{Decomposição do efeito mecânico do preço} sobre o peso e sobre a variação usada no
ajuste parcial, separando-o da decisão ativa do gestor.
\item \textbf{Regra de decisão com o $\lambda$ dentro do fundo} (efeito-fixo de fundo, que se
mostrou estável, e não o Arellano--Bond, que não se mostrou), em vez do agregado, na avaliação de
Diebold--Mariano.
\item \textbf{Diagnosticar a instabilidade do Arellano--Bond} com mais detalhe --- por exemplo,
testando especificações intermediárias de colapso de instrumentos, ou um desenho de
Blundell--Bond (\emph{system GMM}) --- antes de descartar de vez o painel dinâmico como
ferramenta para este problema.
\item \textbf{Modelo em duas partes (\emph{hurdle})} unificando a margem extensiva e a intensiva.
\item \textbf{Sensibilidade do corte de cotistas} ($>1$, $>5$, $>10$, em vez de apenas $>3$).
\end{enumerate}

%==================== REFERÊNCIAS ====================
\newpage
\addcontentsline{toc}{section}{Referências}
\section*{Referências}
\small
\begin{description}
\item Arellano, Manuel and Stephen Bond. Some tests of specification for panel data: Monte Carlo
evidence and an application to employment equations. \emph{Review of Economic Studies},
58(2):277--297, 1991.
\item Cameron, A.\ Colin, Jonah B.\ Gelbach, and Douglas L.\ Miller. Robust inference with
multiway clustering. \emph{Journal of Business \& Economic Statistics}, 29(2):238--249, 2011.
\item Coval, Joshua and Erik Stafford. Asset fire sales (and purchases) in equity markets.
\emph{Journal of Financial Economics}, 86(2):479--512, 2007.
\item Diebold, Francis X.\ and Roberto S.\ Mariano. Comparing predictive accuracy. \emph{Journal
of Business \& Economic Statistics}, 13(3):253--263, 1995.
\item Fama, Eugene F. and James D. MacBeth. Risk, return, and equilibrium: empirical tests.
\emph{Journal of Political Economy}, 81(3):607--636, 1973.
\item Gabaix, Xavier and Ralph S. J. Koijen. In search of the origins of financial fluctuations:
the inelastic markets hypothesis. \emph{NBER Working Paper} 28967, 2021.
\item Harvey, Andrew C. \emph{Forecasting, Structural Time Series Models and the Kalman Filter}.
Cambridge University Press, Cambridge, 1989.
\item Koijen, Ralph S. J. and Motohiro Yogo. A demand system approach to asset pricing.
\emph{Journal of Political Economy}, 127(4):1475--1515, 2019.
\item Lou, Dong. A flow-based explanation for return predictability. \emph{Review of Financial
Studies}, 25(12):3457--3489, 2012.
\item Núcleo de Economia Financeira (NEFIN), USP. Risk factors for the Brazilian market.
\url{https://nefin.com.br/data/risk-factors/}, acessado em 09/07/2026.
\item Newey, Whitney K. and Kenneth D. West. A simple, positive semi-definite, heteroskedasticity
and autocorrelation consistent covariance matrix. \emph{Econometrica}, 55(3):703--708, 1987.
\item Papke, Leslie E. and Jeffrey M. Wooldridge. Econometric methods for fractional response
variables with an application to 401(k) plan participation rates. \emph{Journal of Applied
Econometrics}, 11(6):619--632, 1996.
\end{description}

\end{document}
